% page 624 of the facsimile (image p-623 of the scan)
% block 188.0 x 246.3 mm ; left margin 8.0 mm
\pgc{-12.700mm}{0.000mm}{
\l{\cel{3}616}
\l{}
\l{}
\l{\sou{viburno}. (\sou{bot.)}\cel{2}Arboreto kun folii blanka e beri reda,}
\l{\cel{3}dispozita buketatre, ek la familio \textquotedbl{}kaprifoliacei\textquotedbl{}.}
\l{\cel{3}L.\cel{1}\sou{viburnus}. - EFIS.}
\l{}
\l{vice.\cel{2}(l) qua pre-nominesis por helpar o suplear (sua su-}
\l{\cel{3}perioro nemediata). (2) Indikas \textquotedbl{}remplase di...\textquotedbl{}. -}
\l{\cel{3}(3) (\sou{kande sequata da infinitivo}) indikas ago kontrea}
\l{\cel{3}ye olta quan onu jus parolis, aludis, mencionis.(ex.:}
\l{\cel{3}vice laborar, lu ludis).}
\l{}
\l{\sou{vicero.}\cel{2}(\sou{anat.})\cel{2}Omna organo vivala esencala quan kontenas}
\l{\cel{3}un ek la tri kavaji : kranio, torako od abdomino. - EFIS}
\l{}
\l{\sou{vicio.}\cel{2}Tendenco kustumala a lo etiko-mala. - EFIS.}
\l{}
\l{\sou{vicina}. Qua esas situita proxime. - EFIS.}
\l{\cel{73})}
\l{\sou{vidar}. (\sou{trans}.) Perceptar la imaji quin la lumo-radii qui}
\l{\cel{3}departas de la objekti lumizata formacas sur la fundo di}
\l{\cel{3}la okulo por konvergar adsur retino.- FIR.}
\l{}
\l{\sou{vidamo}. (\sou{feudismo}). Reprezentanto di episkopo, di abado, di}
\l{\cel{3}qua la tasko konsistis ye defensar lu ed, en la kazi maxim}
\l{\cel{3}multa, feudizita da lu. - EFIRS.}
\l{}
\l{\sou{vidva}. Qua cesis havar spozo, pro la morto di ta spozo. -DEFIRS}
\l{}
\l{\sou{vielo}. (\sou{muziko})\cel{2}Muzikilo kordoza, quan onu pleas per klavi}
\l{\cel{3}quin la manuo sinistra igas kontaktar rotaco-roto indu-}
\l{\cel{3}tita ye kolofono qua pleas la rolo di arketo.- DEF.}
\l{}
\l{\sou{vigilar}. (\sou{netrans.}) Ne dormar dum la tempo normale dormo-}
\l{\cel{3}konsakrita. - EFIS.}
\l{}
\l{\sou{vigilio}. (\sou{liturgio katol.}) Hiero di granda festo, dum qua onu}
\l{\cel{3}abstinencas e fastas. - DEFIS.}
\l{}
\l{\sou{vigoro}. Forteso di la korpo, di la spirito o di la paiko, qua}
\l{\cel{3}atingas sua maximo. - EFIS.}
\l{}
\l{\sou{vikario}. (\sou{precipue pri religio}). La persono qua, dum la abs-}
\l{\cel{3}enteso di sua superioro, remplasas lu ed exekutas lua}
\l{\cel{3}taski. - DEFIRS.}
\l{}
\l{\sou{viktimo}. I. Vivanto ofrata sakrifike a la deajo. - II. (me-}
\l{\cel{3}\sou{tafore}). Persono qua sakrifikas vole sua vivo o sua fe-}
\l{\cel{3}liceso ad ulu, ad ulo. - III. Persono qua sakrifikesas a}
\l{\cel{3}la odio, a la venjo, di ulu. - EFIS.}
\l{}
\l{\sou{viktualio}.\cel{2}Singla unajo ek provizuro de nutrivi. - DEFS.}
\l{}
\l{vikuno. (\sou{zool}.) Ruminero sen korni, animalo lanoza ek la ge-}
\l{\cel{3}nero \textquotedbl{}lamao\textquotedbl{}, devenanta de Peru. - L.\cel{1}\sou{auchenia vicunna}}
\l{\cel{3}- DEFIRS.}
}
